% Drawing 67/03-2 -- folder 67, batch 03, archive page 50 (PDF page 51).
% Transcription: au crayon, sur la moitié droite du feuillet (tenu dans le sens de son écriture), trois croquis : en haut une sphère portant les points marqués $1$, $0$, $\infty$ et $\alpha$ (ce dernier cerclé), coupée par un grand cercle presque vertical, dont la moitié arrière est en tirets, et par un équateur ; à gauche, plus bas, un ovale étroit dont un bord est en tirets, dans un contour plus large à peine tracé, marqué $\alpha_0$ ; en bas une sphère dont l'équateur porte $0$, $1$ et, à l'arrière, $\infty$, chacun relié au pôle par un arc.
% Kind: three pencil sketches: two spheres with marked points, an oval. Method 2 (schematic redraw from the facsimile). Status: partial.
% Notes: The sheet is upside down in the scan (his writing runs the other way from the printed listing); the crop is turned by 180 degrees and the sketches are redrawn upright. Top: the sphere (outline as one closed curve; he went over it twice), the equator drawn as its front arc, the nearly vertical great circle (left half full, right half dashed), the points 1 (upper left), 0 (upper right), infinity (on the equator, right of the great circle) and alpha (ringed, at the left end of the equator), a thin line from the point 1 to the great circle carrying two small points, the stroke from alpha to the point 1, an arc on the right from the top through 0 down to the bottom, and two small crosses. Left: the narrow oval (right edge dashed), the wider faint contour round it, and the label alpha_0. Bottom: the sphere, its equator (an ellipse), the pole with a tick above it, the points 0, 1 on the front of the equator and infinity at the back, and the three arcs from the pole. Left out: a fourth, very faint, unlabelled oval at the upper left of the crop; his formulas and the printed listing. Still differs: in the top sphere several faint pencil arcs inside the outline (between alpha, the great circle and the bottom) are not all drawn, and those that are drawn are a reading; outlines are smooth curves through read points. The transcription's note says « le pôle $\infty$ » for the third sphere: on the page infinity sits on the back of the equator and the pole is unlabelled. Redrawn 2026-11-12 (issue 128): to check again.
% Coordinates: 1 unit = 1 mm = 1/100 of the crop width (bbox in page fractions: [0.66, 0.07, 0.96, 0.52]).
% Redrawn by Claude (claude-fable-5-2), issue #128, 2026-11-12; not checked by a human; not compiled (no TeX in the session).
% Rights: reproduces part of the work -- behind RENDER_TRANSCRIPTS, never published by itself.
\documentclass[tikz,border=3pt]{standalone}
\usepackage{amsmath,amssymb}
\usetikzlibrary{patterns,arrows.meta}
\tikzset{pencil/.style={black!45}, struck/.style={black!35}, mpt/.style={circle,draw,inner sep=1.3pt}, pt/.style={circle,fill,inner sep=0.9pt}, opt/.style={circle,draw,fill=white,inner sep=1.1pt}}
\begin{document}
\begin{tikzpicture}[x=1mm,y=1mm,line cap=round,line join=round,font=\scriptsize]
\draw[pencil] plot[smooth cycle, tension=0.55] coordinates {(70,-5) (80,-8) (87,-16) (90,-27) (88,-38) (82,-46) (72,-50) (62,-49) (53,-42) (48,-32) (48,-22) (53,-12) (61,-6)};
\draw[pencil] plot[smooth, tension=0.55] coordinates {(47.5,-26) (50,-30) (63.5,-31) (76.5,-31) (89,-30.5)};
\draw plot[smooth, tension=0.55] coordinates {(72,-5.5) (66,-15) (63.5,-31) (64,-42) (66.5,-49.5)};
\draw[dashed] plot[smooth, tension=0.55] coordinates {(72,-5.5) (74,-14) (73,-31) (70,-43) (66.5,-49.5)};
\draw[pencil] (58.5,-22.5) -- (68,-21.5) -- (73,-22);
\node[pt] at (68,-21.5) {};
\node[pt] at (73,-22) {};
\draw[pencil] (52.5,-30.5) -- (58.5,-22.5);
\draw[pencil] plot[smooth, tension=0.55] coordinates {(80,-9) (81.5,-20) (82.5,-32) (79,-43) (74,-49)};
\node[pt] at (58.5,-22.5) {};
\node[inner sep=1pt] at (58,-17.5) {\small $1$};
\node[pt] at (83,-23.5) {};
\node[inner sep=1pt] at (86.5,-22) {\small $0$};
\node[pt] at (76.5,-31) {};
\node[inner sep=1pt] at (76.5,-34.5) {\small $\infty$};
\node[pt] at (63.5,-31) {};
\node[mpt] at (52.5,-30.5) {};
\node[inner sep=1pt] at (52.5,-34) {\small $\alpha$};
\draw (81,-13) -- (83,-13);
\draw (82,-12) -- (82,-14);
\draw (76,-44.5) -- (78,-44.5);
\draw (77,-43.5) -- (77,-45.5);
\draw plot[smooth, tension=0.55] coordinates {(33,-114) (26,-120) (22,-135) (21.5,-150) (24,-161) (27,-163.5)};
\draw[dashed] plot[smooth, tension=0.55] coordinates {(33,-114) (38,-125) (39,-140) (36,-155) (27,-163.5)};
\draw[pencil] plot[smooth, tension=0.55] coordinates {(33,-114) (18,-120) (8,-133) (6,-145) (10,-158) (20,-165) (27,-164)};
\draw[pencil] plot[smooth, tension=0.55] coordinates {(33,-114) (44,-118) (52,-130) (53,-141) (48,-151)};
\node[inner sep=1pt] at (5.5,-137.5) {\small $\alpha_0$};
\draw[pencil] plot[smooth cycle, tension=0.55] coordinates {(68.5,-154) (80,-157) (88,-166) (91,-178) (88,-191) (80,-200) (68,-203.5) (56,-201) (46,-193) (42,-180) (45,-167) (54,-158)};
\draw[pencil] (66.5,-178.5) ellipse[x radius=23, y radius=7.2];
\node[pt] at (68.5,-154) {};
\draw (67.2,-152) -- (68.3,-151) -- (69.6,-152);
\node[pt] at (50,-182.5) {};
\node[inner sep=1pt] at (52,-186.5) {\small $0$};
\node[pt] at (83,-183) {};
\node[inner sep=1pt] at (83,-188) {\small $1$};
\node[pt] at (68.5,-171.3) {};
\node[inner sep=1pt] at (73.5,-169) {\small $\infty$};
\draw plot[smooth, tension=0.55] coordinates {(68.5,-154) (58,-160) (52.5,-170) (50,-182.5)};
\draw plot[smooth, tension=0.55] coordinates {(68.5,-154) (70.5,-160) (70,-166) (68.5,-171.3)};
\draw plot[smooth, tension=0.55] coordinates {(68.5,-154) (79,-158) (83.5,-168) (83,-183)};
\end{tikzpicture}
\end{document}
